\section{Case Study}

The failure analysis compares Enhanced Hybrid w122 with Enhanced KCH-MedRank at top 10 over 943 paired held-out BioASQ questions. KCH-MedRank rescues 661 gold passages into top 10 and loses 307 gold passages from top 10. This net gain supports the retrieval claim, but the lost cases are important because they show where knowledge-constrained features can overrule a strong lexical or semantic rank.

\begin{table}[t]
\centering
\small
\caption{Top-10 failure analysis comparing Enhanced Hybrid w122 with Enhanced KCH-MedRank on 943 paired held-out questions.}
\label{tab:failure-analysis-summary}
\begin{tabular}{lr}
\toprule
Diagnostic & Value \\
\midrule
Paired questions & 943 \\
Questions with rescued gold evidence & 414 \\
Rescued gold passages & 661 \\
Questions with lost gold evidence & 214 \\
Lost gold passages & 307 \\
Net rescued passages & +354 \\
\bottomrule
\end{tabular}
\end{table}


\begin{figure}[t]
\centering
\resizebox{\linewidth}{!}{%
\begin{tikzpicture}[
  font=\small,
  panel/.style={draw=black!18, rounded corners=3pt, fill=black!2},
  chip/.style={draw=black!25, rounded corners=2pt, fill=white, align=center, inner xsep=5pt, inner ysep=3pt},
  good/.style={draw=green!55!black, fill=green!12},
  bad/.style={draw=orange!70!black, fill=orange!12},
  arrow/.style={-{Latex[length=2mm]}, line width=0.75pt},
  lab/.style={font=\footnotesize, text=black!70}
]
\node[panel, minimum width=6.9cm, minimum height=4.2cm, anchor=south west] at (0,0) {};
\node[anchor=west, font=\bfseries] at (0.25,3.82) {Rescued evidence: p-crAssphage};
\node[chip] (h54) at (1.15,2.55) {Hybrid\\rank 54};
\node[chip, good] (k1) at (5.75,2.55) {KCH-MedRank\\rank 1};
\draw[arrow, draw=green!50!black] (h54) -- node[above, lab] {promoted into top 10} (k1);
\node[chip, good] at (1.25,1.25) {bacteriophage};
\node[chip, good] at (3.15,1.25) {crAssphage};
\node[chip, good] at (5.0,1.25) {viral genome};
\node[lab, align=center] at (3.45,0.48) {Shared biomedical concepts and a local phage-related candidate cluster\\support the evidence promotion.};

\node[panel, minimum width=6.9cm, minimum height=4.2cm, anchor=south west] at (7.35,0) {};
\node[anchor=west, font=\bfseries] at (7.6,3.82) {Lost evidence: ubiquitination};
\node[chip, bad] (h2) at (8.45,2.55) {Hybrid\\rank 2};
\node[chip] (k18) at (13.05,2.55) {KCH-MedRank\\rank 18};
\draw[arrow, draw=orange!75!black] (h2) -- node[above, lab] {demoted outside top 10} (k18);
\node[chip, bad] at (8.65,1.25) {mono-/poly-\\ubiquitination};
\node[chip, bad] at (10.85,1.25) {ubiquitin};
\node[chip, bad] at (12.85,1.25) {post-translational\\processing};
\node[lab, align=center] at (10.85,0.48) {Surface-form mismatch and incomplete concept normalization\\weaken the knowledge-alignment signal.};
\end{tikzpicture}}
\caption{Visual case study of one rescued and one lost BioASQ evidence passage. The figure summarizes how local biomedical concept alignment can promote a missed gold passage, while incomplete synonym and hierarchy alignment can also demote a relevant passage.}
\label{fig:case-study-path}
\end{figure}


\begin{table}[t]
\centering
\small
\caption{Representative rescued and lost evidence cases from the held-out BioASQ comparison between Enhanced Hybrid w122 and KCH-MedRank.}
\label{tab:case-study}
\begin{tabularx}{\linewidth}{p{0.11\linewidth}p{0.29\linewidth}p{0.19\linewidth}p{0.14\linewidth}X}
\toprule
Case & Question & Key entities or MeSH signals & Gold rank & Interpretation \\
\midrule
Rescued & What is the p-crAssphage? & bacteriophages; crAssphage; viral genome; phylogeny & 54 $\rightarrow$ 1 & The gold passage shares the bacteriophage concept with the question and is connected to a local cluster of phage-related candidates, allowing the reranker to recover evidence missed by the enhanced hybrid ordering. \\
Lost & Is there any protein that undergoes both mono-ubiquitination and poly-ubiquitination? & ubiquitin; ubiquitination; ubiquitins; post-translational processing & 2 $\rightarrow$ 18 & The passage is relevant but uses related biochemical forms and broader processing concepts; synonym and hierarchy features do not fully align mono-/poly-ubiquitination with the passage evidence. \\
\bottomrule
\end{tabularx}
\end{table}


Figure~\ref{fig:case-study-path} and Table~\ref{tab:case-study} illustrate one rescued case and one lost case. In the rescued p-crAssphage example, the gold passage moves from rank 54 to rank 1. The main alignment signal is the shared bacteriophage concept, supported by a local cluster of phage-related candidates. This pattern is consistent with the method design: MeSH and local hypergraph features can promote evidence that is medically aligned even when the initial enhanced hybrid ordering places it outside the top 10.

The lost ubiquitination case shows the opposite behavior. A gold passage drops from rank 2 to rank 18 even though it discusses mono-ubiquitinated LDH-A and poly-ubiquitinated LDH. The question and passage contain related biochemical terminology, but the synonym and hierarchy features do not fully align the surface forms and broader post-translational-processing concepts. Similar losses occur for thyroid hormone analogs, cilengitide receptors, dasatinib treatment, GIST treatment, and Fanconi anemia. These examples motivate future improvements in biomedical concept normalization rather than stronger claims about the relation features.
